\chapter{Git and GitHub}
\label{ch:git}

\chapterbrief
  {Put your work under version control and back it up.}
  {Commits as save points; SSH keys, a second time.}
  {\code{git push} working with no password.}

\section{Why bother, on a solo project}

Three reasons, and only the first is the obvious one.

\begin{description}[leftmargin=0pt, style=nextline]
  \item[Backup.] SD cards die. Laptops get spilled on.
  \item[Undo.] Not \code{Ctrl-Z} undo --- \emph{``the reader worked on Tuesday
        and does not now''} undo. Git can show you exactly what changed between
        Tuesday and today.
  \item[A record of your reasoning.] In Volume One each chapter ends with a
        commit. Six months later, \code{git log} is a narrative of what you built
        and in what order --- which is worth more than you expect when you come
        back to it, or film it.
\end{description}

\begin{analogy}{Save points, not autosave}
A word processor autosaves continuously; you cannot go back to ``before I
rewrote that section''.

Git is save points in a game. \emph{You} decide when a moment is worth marking,
and you name it. ``M2: driver reads the version register'' is a save point you
can return to, compare against, or branch from.

The skill is choosing when to save: after something works, before something
risky.
\end{analogy}

\section{Start the repository}

You will do this inside the project folder from Chapter~\ref{ch:project}. If you
have not created it yet, come back after that chapter.

\begin{shell}
# on your laptop, inside the project folder
git init
git status
\end{shell}

\code{git status} is the command you will run most. It answers: what have I
changed, and what is staged?

\section{Your first commit}

\begin{shell}
git add .                 # stage everything not ignored
git status                # confirm what is staged -- always look
git commit -m "chore: project skeleton from the starter pack"
git log --oneline
\end{shell}

\begin{concept}{Why two steps, \texttt{add} then \texttt{commit}}
\code{add} stages changes; \code{commit} records what is staged.

The separation exists so you can commit \emph{part} of your work. You fixed a
bug and also renamed a variable; those are two different stories, and staging
lets you tell them separately.

At first, \code{git add .} then \code{git commit} is fine. Later, \code{git add
-p} steps through each change and asks whether to include it --- and that is when
the two-step design starts earning its keep.
\end{concept}

\subsection*{Commit messages that help later}

A convention worth adopting, because it makes \code{git log} scannable:

\begin{trace}
feat:  a new capability          feat: read the card UID
fix:   a bug fixed               fix: skip the 0x88 cascade tag
chore: housekeeping              chore: add .gitignore
docs:  documentation             docs: explain the FIFO wait loop
test:  tests only                test: RFC 4493 vectors for CMAC
\end{trace}

Write the message as the completion of \emph{``this commit will\ldots''} ---
present tense, imperative. \code{fix: skip the cascade tag}, not \code{fixed the
cascade tag}.

\begin{trap}{What must never be committed}
\code{.gitignore} in the starter pack already excludes \code{target/} and common
secret filenames. Two habits on top of that:

\begin{itemize}
  \item \textbf{Never commit a private key.} \code{id\_ed25519} stays in
        \code{\textasciitilde/.ssh/}, never in a project.
  \item \textbf{Never commit a password.} Not in a config file, not in a
        comment, not ``temporarily''.
\end{itemize}

Git history is permanent. Deleting a secret in a later commit does \emph{not}
remove it --- it is still there in the history, and if you pushed it publicly,
assume it is compromised and rotate it.
\end{trap}

\section{GitHub}

Create an account at \code{github.com} if you do not have one, then a new
repository: \textbf{New} $\rightarrow$ name it \code{midnight-brewer}
$\rightarrow$ \textbf{Private} $\rightarrow$ and \textbf{do not} tick any of
``add a README / .gitignore / licence'' (you already have files; those options
create a conflicting first commit).

\subsection{An SSH key for GitHub}

You can push over HTTPS with a token, but SSH is less friction and you already
know how it works.

\begin{shell}
# on your laptop -- reuse the key you already made
cat ~/.ssh/id_ed25519.pub
\end{shell}

Copy the whole line (it starts \code{ssh-ed25519} and ends with your comment).
In GitHub: \textbf{Settings} $\rightarrow$ \textbf{SSH and GPG keys}
$\rightarrow$ \textbf{New SSH key}, paste, save.

Test:

\begin{shell}
ssh -T git@github.com
\end{shell}

\begin{trace}
Hi evilcomputer12! You've successfully authenticated,
but GitHub does not provide shell access.
\end{trace}

That message is success --- GitHub is not supposed to give you a shell.

\begin{concept}{The same padlock, a second lock}
This is the identical key pair from Chapter~\ref{ch:sshkeys}. You handed a copy
of the public half to your Pi; now you have handed another copy to GitHub.

One private key, many places that recognise it. That is the whole model, and it
is why the private key never moving matters: it is the single thing whose
compromise would unlock everything at once.
\end{concept}

\section{Push}

\begin{shell}
git remote add origin git@github.com:YOURNAME/midnight-brewer.git
git branch -M main
git push -u origin main
\end{shell}

\begin{trace}
Enumerating objects: 14, done.
Writing objects: 100% (14/14), 3.21 KiB | 3.21 MiB/s, done.
To github.com:YOURNAME/midnight-brewer.git
 * [new branch]      main -> main
branch 'main' set up to track 'origin/main'.
\end{trace}

Refresh the GitHub page --- your files are there.

The \code{-u} sets \code{origin/main} as the default, so from now on it is just:

\begin{shell}
git push
\end{shell}

\section{The loop you will actually use}

\begin{shell}
git status                    # what changed?
git diff                      # show me the actual changes
git add -A                    # stage everything
git commit -m "feat: ..."     # save point
git push                      # back it up
\end{shell}

And two that save you:

\begin{shell}
git log --oneline -10         # last ten save points
git diff HEAD~1               # what changed in the last commit
git checkout -- file.java     # throw away my changes to this file
\end{shell}

\begin{trap}{\texttt{git checkout -{}-} is destructive}
That last command discards your uncommitted changes to that file,
\emph{permanently}. There is no undo, because the changes were never recorded.

That is precisely why you commit often: committed work can always be recovered;
uncommitted work is one careless command from gone.
\end{trap}

\begin{checkpoint}{Your code is backed up}
\begin{shell}
git log --oneline
git push
\end{shell}

At least one commit exists, \code{push} reports \emph{Everything up-to-date},
and you can see your files on github.com.
\end{checkpoint}
